\part{Clôture dodécaphasique et dépendances transversales de l'état}

\chapter[Dérivation interne de la constante de structure fine]{Dérivation interne de \texorpdfstring{\(\alpha_{\rm HMT}\)}{alpha HMT} : généalogie de survie, clôture dodécaphasique, deux constructions compatibles et bifurcation de Witt}
\label{chap:alpha-rev7}
\HMTClaim{HMT-R-ALPHA-DODECAPHASE}
\HMTClaim{HMT-R-DODEC-EXTRACTOR}

\paragraph{Quatre strates de \(\alpha_{\rm HMT}\).}
\begin{enumerate}
 \item Le constructeur généalogique est l'état terminal enrichi muni de ses deux expressions parallèles \(B_{\rm term}\) et \(U_{12}^{\rm sgn}\).
 \item La caractérisation analytique est la chaîne réversible \(U_{12}^{\rm sgn}\leftrightarrow K\to A\to\alpha_{12}\), suivie de son prolongement coinductif.
 \item Étant adimensionnelle, sa transduction commune n'introduit pas d'unité : elle fixe la carte angulaire \(A_{\rm deg}=1000\alpha_{\rm HMT}\) utilisée par les canaux ultérieurs.
 \item La comparaison avec CODATA ne s'ouvre qu'après l'obtention de \(\alpha_{\rm HMT}\) ; elle ne sélectionne ni \(K\), ni la retenue, ni le prolongement.
\end{enumerate}
La généalogie complète conserve explicitement la chaîne d'entrée et atteint un état terminal enrichi. Ses deux expressions sont parallèles :
\[
\begin{aligned}
 w_6\longrightarrow w_{30}\longrightarrow R_{36}\longrightarrow G_9
 \longrightarrow x_{\rm term}
 &\xrightarrow{\ \pi_{\rm term}\ }B_{\rm term},\\
 x_{\rm term}
 &\xrightarrow{\ R_{\rm sgn}\ }U_{12}^{\rm sgn}
 \xleftrightarrow{\ \mathcal H_{12}\ }K
 \longrightarrow A\longrightarrow\alpha_{12}.
\end{aligned}
\]
\(B_{\rm term}\) n'engendre pas \(U_{12}^{\rm sgn}\) : tous deux conservent des coordonnées différentes du même état. La transformation \(\mathcal H_{12}\) est inversible, et l'extracteur \((D_3K,D_4K,Q(K))\) est de rang \(12\). Dans ce domaine de primitives terminales explicites, la chaîne ne reçoit en entrée ni \(\alpha\), ni CODATA, ni une masse, ni aucune autre coordonnée métrologique. L'évaluation archimédienne forme \((P,E,\Phi,K,c_{13}=0)\), applique le cocycle de retenue et dérive \(\alpha_{\rm HMT}\). L'évaluation d'incidence transporte les mêmes mots et le même sceau vers \(B_K\subset H_\alpha^-\), qui constituera l'entrée de la branche Witt--Golay. La construction E21/22 et le prolongement distant \(\mathsf T_{\alpha,\infty}\) ---appelé \(T1\) dans la notation originale--- conservent leurs propres domaines et applications comme développements de la généalogie dodécaphasique principale.

\HMTDemoteHeadings
\HMTInputLive{sections/hmt/08_dodecafase_alpha}
\HMTInputLive{sections/hmt/08a_alpha_dos_vias_t1_rev8}
\HMTInputLive{sections/hmt/08b_alpha_t1_10000_rev9}
\HMTRestoreHeadings

\HMTInputLive{sections/ampliacion_20260818/03_extension_holografica_euler}

\HMTInputLive{sections/hmt/07b_realizacion_analitica_contraangulo}
\HMTInputLive{sections/hmt/16c_transportador_positivo_accion}
\HMTInputLive{sections/hmt/delta_297/16c_entrelazador_dodecafase_witt}
\begingroup
\def\HMTInternalTraceability{1}
\HMTInputLive{sections/hmt/delta_297/16d_medida_accion_positiva}
\endgroup

\HMTInputLive{sections/md/04b_banda_particula_virtual}

\chapter*{Construction holographique complète de l'électron comme dépendance du lecteur}
\addcontentsline{toc}{chapter}{Construction holographique complète de l'électron comme dépendance du lecteur}
\HMTInputLive{sections/md/06f_biografia_electron_rev6}

\chapter*{Transduction métrologique des états généalogiques}
\addcontentsline{toc}{chapter}{Transduction métrologique des états généalogiques}
\HMTInputLive{sections/ampliacion_20260818/07_transduccion_metrologica_estados_genealogicos}
